\section{Data and Methodology}\label{sec:data_method}

\subsection{Universe, sample and data sources}

Raw daily data are obtained from Yahoo Finance over the period 2010-01-05 to
2026-03-30. The initial universe comprises 78 equity tickers spanning 20 European
countries, drawn from benchmark indices (EURO STOXX Banks, STOXX Europe 600 Financials)
supplemented with major national institutions. We apply a set of filters to construct a balanced weekly panel. An institution~$i$ is
retained if (i)~its price series starts before a minimum calendar date (configured as
\texttt{min\_start = 2010-01-01}), (ii)~the weekly series obtained after preprocessing
contains at least $\texttt{min\_obs} = 252$ observations, (iii)~the fraction of missing
weekly observations does not exceed $\texttt{drop\_na\_ratio\_max} = 0.20$, and
(iv)~the stock is not subject to prolonged suspensions or structural breaks (such as
delisting) within the estimation period.%
\footnote{The parameter \texttt{remove\_late = TRUE} ensures that late entrants are
dropped rather than padded, and \texttt{keep\_min\_symbols = 3} prevents countries with
fewer than three surviving institutions from entering the analysis.}
Three tickers were excluded due to data quality issues (SVEG.OL, VMUK.L, TRIO.AS), and
two were unavailable from the data provider at retrieval time (BOV.MT, SYDB.CO). After
all filters are applied, the final panel retains 63 institutions from 18 countries,
observed over 682 weekly dates (December 2012 -- March 2026 after the first valid
return is computed). Survivorship bias is partially mitigated by including institutions
that were subsequently acquired or delisted within the sample period, provided they
satisfy the inclusion criteria above. The full list of institutions is provided in
Appendix~\ref{app:sample}, and their descriptive return statistics are reported in
Appendix~\ref{app:sumstats}.

Table~\ref{tab:panel_summary} summarises the cross-sectional distribution of return
properties. The average weekly standard deviation of 4.50\% reflects the high intrinsic
volatility of European financial equities. Returns are, on average, negatively skewed
($-0.83$) and exhibit pronounced excess kurtosis ($8.39$), consistent with fat-tailed
distributions driven by crisis episodes. The cross-sectional range of standard deviations
(1.42\%--11.92\%) highlights the heterogeneity between stable cantonal banks and highly
volatile peripheral institutions.

% Auto-generated by 16_summary_statistics.R
\begin{table}[htbp]
  \centering
  \caption{Panel summary of weekly log-returns. Statistics are computed
  over the full sample period and then averaged across institutions.
  $N=63$ institutions, 18 countries, weekly frequency.}
  \label{tab:panel_summary}
  \begin{tabular}{lrr}
  \hline
  \textbf{Statistic} & \textbf{Cross-section mean} & \textbf{Cross-section range} \\
  \hline
  Observations per series & 682 & [682, 682] \\
  Mean return (\%) & 0.1347 & [-1.1695, 0.4249] \\
  Std deviation (\%) & 4.504 & [1.423, 11.916] \\
  Minimum weekly return (\%) & -32.38 & [-115.82, -5.78] \\
  Maximum weekly return (\%) & 21.45 & [6.41, 84.73] \\
  Skewness & -0.827 & [-4.075, 1.413] \\
  Excess kurtosis & 8.39 & [1.41, 47.98] \\
  \hline
  \end{tabular}
\end{table}


For each institution $i \in \{1,\dots,N\}$ and trading day $\tau$, we collect adjusted
closing prices $P_{i,\tau}$ and, when available, high--low--open--close (HLOC) prices
$(O_{i,\tau}, H_{i,\tau}, L_{i,\tau}, C_{i,\tau})$. The raw series are checked for
missing observations and obvious outliers; when necessary, isolated missing values are
linearly interpolated or removed, and corporate actions (splits, large dividends) are
controlled for when not already embedded in adjusted prices. Daily trading days are then
aligned across institutions and aggregated to a common weekly calendar, so that the final
time index $t = 1,\dots,T$ refers to consecutive weeks between $\tau_0$ and $\tau_1$
(2010--2026 in our application). All institutions are therefore observed on the same set
of weekly dates, which avoids spurious jumps in co-movement arising from asynchronous
trading or late entries.

\subsection{Returns and volatility construction}

For each institution $i$ and trading day $\tau$, we define the daily log-return based on
adjusted closing prices as
\begin{equation}
  r_{i,\tau}
  = \log P_{i,\tau} - \log P_{i,\tau-1},
  \qquad i = 1,\dots,N,\quad \tau = 2,\dots,T^d.
\end{equation}
Weekly returns are constructed by taking the Friday-to-Friday log-difference (or the last
trading day of each week). These weekly log-returns constitute the input series used to
estimate connectedness.

Unlike studies that rely on range-based (Parkinson) or realized-volatility estimators, we
deliberately do \emph{not} compute volatility proxies from high--low or intraday data.
Following standard practice in VAR-based spillover analysis, we model the dynamics of
weekly returns directly: when placed inside a rolling VAR, their conditional second
moments serve as volatility proxies and allow us to capture the propagation of
uncertainty across institutions.

Let $\mathbf{r}_t = (r_{1,t},\dots,r_{N,t})'$ denote the vector of weekly returns. For
numerical stability and comparability across institutions, each series is standardized to
have zero mean and unit variance over the sample:
\[
  y_{i,t} = \frac{r_{i,t} - \bar{r}_i}{s_i},
\]
where $\bar{r}_i$ and $s_i$ denote the sample mean and standard deviation of institution
$i$. The VAR--DY connectedness measures are computed on the standardized return vector
$\mathbf{y}_t = (y_{1,t},\dots,y_{N,t})'$.

\subsection{\texorpdfstring{VAR($p$)}{VAR(p)} specification and generalized FEVD}\label{subsec:var_gfevd}

The connectedness measures are derived from a vector autoregression of order $p$ for the
standardized weekly return vector $\mathbf{y}_t$, which we interpret as a return-based
volatility proxy.

Formally, on a probability space $(\Omega,\mathcal{F},\mathbb{P})$ with filtration
$(\mathcal{F}_t)_{t\geq 0}$ generated by $\{\mathbf{y}_s : s \le t\}$, we assume that
\begin{equation}
  \mathbf{y}_t
  = \sum_{k=1}^p A_k \mathbf{y}_{t-k} + \boldsymbol{\varepsilon}_t,
  \qquad t = p+1,\dots,T,
  \label{eq:varp}
\end{equation}
where $A_k$ are $N \times N$ coefficient matrices and $\boldsymbol{\varepsilon}_t$ is an
$N$-dimensional innovation process satisfying
\begin{align}
  \mathbb{E}[\boldsymbol{\varepsilon}_t \mid \mathcal{F}_{t-1}] &= \mathbf{0}, \\
  \mathbb{E}[\boldsymbol{\varepsilon}_t \boldsymbol{\varepsilon}_t' \mid \mathcal{F}_{t-1}] &= \Sigma, \\
  \mathbb{E}[\boldsymbol{\varepsilon}_t \boldsymbol{\varepsilon}_s'] &= \mathbf{0}
    \quad \text{for } t \neq s,
\end{align}
with $\Sigma$ a positive-definite $N \times N$ covariance matrix. In other words,
$\{\boldsymbol{\varepsilon}_t\}$ is a zero-mean, covariance-stationary weekly white-noise
process with finite second moments. For likelihood-based inference one may additionally
assume conditional normality, $\boldsymbol{\varepsilon}_t \mid \mathcal{F}_{t-1} \sim
\mathcal{N}(0,\Sigma)$, but the GFEVD-based connectedness measures do not require
Gaussianity.

In the empirical implementation we set $p = 1$. Lag order selection evaluates AIC and BIC
over $p \in \{1,\dots,8\}$ on the full standardized returns panel; both criteria
consistently select $p = 1$. Fixing $p = 1$ also improves numerical stability in rolling
estimation.%
\footnote{Replicating the analysis with $p = 2$ yields virtually identical TCI dynamics
and institutional rankings.}
The lag order satisfies the stability requirement that the eigenvalues of the companion
matrix lie strictly inside the unit circle. Under this stability condition, the VAR admits a
moving-average representation
\begin{equation}
  \mathbf{y}_t = \sum_{h=0}^{\infty} \Phi_h \boldsymbol{\varepsilon}_{t-h},
  \qquad \Phi_0 = I_N,
\end{equation}
with $\Phi_h$ obtained recursively from the $\{A_k\}_{k=1}^p$.

Following \citet{DieboldYilmaz2012}, we employ the Generalized Forecast Error Variance
Decomposition (GFEVD) of \citet{KoopPesaranShin1996} and \citet{PesaranShin1998}, which
is invariant to the ordering of variables. Let $\sigma_{jj}$ denote the $j$th diagonal
element of $\Sigma$ and $\mathbf{e}_j$ the $j$th column of the $N \times N$ identity
matrix. The $H$-step-ahead generalized variance share from variable $j$ to variable $i$
is defined as
\begin{equation}
  \theta_{ij}^{g}(H)
  = \frac{
      \sigma_{jj}^{-1}
      \sum_{h=0}^{H-1}
      \left( \mathbf{e}_i' \Phi_h \Sigma \mathbf{e}_j \right)^2
    }{
      \sum_{h=0}^{H-1}
      \left( \mathbf{e}_i' \Phi_h \Sigma \Phi_h' \mathbf{e}_i \right)
    },
  \qquad i,j = 1,\dots,N.
\end{equation}
Because the generalized shares for a given $i$ do not necessarily sum to one across $j$,
we normalize row-wise:
\begin{equation}
  \tilde{\theta}_{ij}^{g}(H)
  = \frac{\theta_{ij}^{g}(H)}
         {\sum_{j=1}^N \theta_{ij}^{g}(H)},
  \qquad \sum_{j=1}^N \tilde{\theta}_{ij}^{g}(H) = 1.
\end{equation}
The normalized GFEVD matrix at horizon $H$ is denoted
\[
  \tilde{\Theta}^g(H)
  = \big[ \tilde{\theta}_{ij}^{g}(H) \big]_{i,j=1}^N.
\]

\noindent\textit{Interpretation caveat.} GFEVD-based connectedness measures predictive
co-movement, defined as the share of $H$-step-ahead forecast error variance of institution $i$
attributable to shocks in institution $j$, and does not identify structural causal
relationships. The directionality in FROM/TO spillovers reflects Granger-causal
predictability, subject to the usual omitted-variable and model-specification caveats
inherent in reduced-form VAR analysis.

\subsection{Connectedness measures and rolling windows}\label{subsec:rolling}

From $\tilde{\Theta}^g(H)$ we construct the standard Diebold--Yilmaz connectedness
indices. The Total Connectedness Index (TCI) is defined as
\begin{equation}
  \mathrm{TCI}(H)
  = \frac{1}{N}
    \sum_{i=1}^N \sum_{\substack{j=1 \\ j \neq i}}^N
    \tilde{\theta}_{ij}^g(H),
\end{equation}
which is the average share of forecast error variance arising from cross-variable shocks
rather than own shocks.

Directional connectedness \emph{from} all others \emph{to} institution $i$ is
\begin{equation}
  C_{i \leftarrow \bullet}(H)
  = \sum_{\substack{j=1 \\ j \neq i}}^N \tilde{\theta}_{ij}^g(H),
\end{equation}
while directional connectedness \emph{to} all others \emph{from} $i$ is
\begin{equation}
  C_{i \rightarrow \bullet}(H)
  = \sum_{\substack{j=1 \\ j \neq i}}^N \tilde{\theta}_{ji}^g(H).
\end{equation}
The net spillover of institution $i$ to the system is then
\begin{equation}
  C_i^{\text{net}}(H)
  = C_{i \rightarrow \bullet}(H) - C_{i \leftarrow \bullet}(H).
\end{equation}

To capture time variation in spillovers, we estimate the VAR model~\eqref{eq:varp} on
rolling windows of fixed length $W$ weekly observations. In the baseline configuration we
set $W = 156$ weeks (approximately three years of data), consistent with \citet{DieboldYilmaz2014}
and prior rolling-VAR connectedness studies; this window length provides sufficient
observations for reliable GFEVD estimation while preserving meaningful time variation.%
\footnote{Robustness to alternative window sizes $W \in \{104, 208\}$ weeks is reported
in Section~\ref{sec:results}; TCI dynamics and institutional rankings are preserved.}
The rolling step is four weeks between consecutive windows.%
\footnote{In practice, connectedness is computed at every weekly observation and then
reported at four-week intervals to reduce output volume; results are numerically identical
to directly stepping by four weeks.}
Let the first window cover $t = 1,\dots,W$, the second window $t = 1+4,\dots,W+4$, and
so on, up to the last window ending at $T$. This yields $W^\ast$ overlapping windows
indexed by $w = 1,\dots,W^\ast$. For each window $w$ we estimate the VAR, compute the
normalized GFEVD matrix $\tilde{\Theta}^g_w(H)$ at forecast horizon $H = 10$ weeks (the
horizon over which the generalized forecast error variance is decomposed; following
\citet{DieboldYilmaz2012}, results are robust to alternative choices $H \in \{4,6,12\}$),
and obtain the corresponding connectedness measures
\[
  \mathrm{TCI}_w,\quad
  C_{i \leftarrow \bullet, w},\quad
  C_{i \rightarrow \bullet, w},\quad
  C_{i,w}^{\text{net}},\qquad i = 1,\dots,N.
\]
These rolling-window indices form time series that summarize the evolution of system-wide
and institution-specific spillovers and serve as inputs to the clustering and comparison
analyses.

\subsection{DCC--GARCH benchmark}\label{subsec:dcc}

As a covariance-based benchmark, we estimate a multivariate DCC--GARCH model on rolling
windows of the weekly return vector $\mathbf{r}_t$. For each rolling window, returns are
standardized using the window-specific mean and standard deviation before DCC--GARCH
estimation, consistent with standard practice for rolling GARCH models (preventing
look-ahead bias from full-sample moments). The DCC model is specified as follows.

Each return series is modeled by a univariate GARCH(1,1):
\begin{equation}
  r_{i,t} = \mu_i + \varepsilon_{i,t}, \qquad
  \varepsilon_{i,t} = h_{i,t}^{1/2} z_{i,t},
\end{equation}
with conditional variance dynamics
\begin{equation}
  h_{i,t}
  = \omega_i + \alpha_i \varepsilon_{i,t-1}^2 + \beta_i h_{i,t-1},
  \qquad \omega_i > 0,\ \alpha_i,\beta_i \ge 0,\ \alpha_i + \beta_i < 1,
\end{equation}
and standardized innovations $z_{i,t}$ such that
\[
  \mathbb{E}[z_{i,t} \mid \mathcal{F}_{t-1}] = 0, \qquad
  \mathbb{E}[z_{i,t}^2 \mid \mathcal{F}_{t-1}] = 1,
\]
assumed conditionally i.i.d.\ Normal (with Student-$t$ as fallback when Normal
initialization fails to converge). Collecting conditional standard deviations in
\[
  D_t = \mathrm{diag}\big(h_{1,t}^{1/2},\dots,h_{N,t}^{1/2}\big),
\]
we write the conditional covariance matrix as
\begin{equation}
  H_t = D_t R_t D_t,
\end{equation}
where $R_t$ is the $N \times N$ conditional correlation matrix.

DCC dynamics \citep{Engle2002DCC} specify an auxiliary process $Q_t$ as
\begin{equation}
  Q_t = (1 - a - b)\bar{Q} + a (z_{t-1} z_{t-1}') + b Q_{t-1},
\end{equation}
where $z_t = D_t^{-1} \boldsymbol{\varepsilon}_t$ is the vector of standardized
residuals, $\bar{Q}$ is the unconditional covariance matrix of $z_t$, and $a,b \ge 0$
satisfy $a + b < 1$ to ensure stationarity. The correlation matrix is obtained by
rescaling:
\begin{equation}
  R_t = \mathrm{diag}(Q_t)^{-1/2} Q_t \,\mathrm{diag}(Q_t)^{-1/2}.
\end{equation}
Compared with BEKK-type specifications \citep{EngleKroner1995}, DCC introduces only two
additional scalar parameters $(a,b)$, independent of $N$, which is crucial in a
high-dimensional setting with rolling estimation.

GARCH and DCC parameters are jointly estimated using the \texttt{rmgarch} package in
\textsf{R}. When the default solver fails to converge, the estimation routine
sequentially tries alternative solver configurations (\texttt{solnp}, \texttt{nlminb},
hybrid combinations) and, if necessary, switches the distributional assumption from
Normal to Student-$t$. Windows for which all solver combinations fail to converge are
excluded from the comparison analysis.

For the comparison with VAR--DY connectedness, we align the DCC output with the VAR
windows by computing, for each window $w$, the mean absolute off-diagonal correlation
from $\{R_t\}$ over the same calendar weeks used in the VAR estimation window. This
yields a sequence $\{\overline{\rho}_w\}_{w=1}^{W^\ast}$ that can be directly compared
to $\{\mathrm{TCI}_w\}$.

\subsection{Correlation-based networks and minimum spanning trees}\label{subsec:mst}

To obtain a purely statistical, model-free view of dependence, we construct
correlation-based networks and extract their backbone topology through minimum spanning
trees (MSTs) in the spirit of \citet{Mantegna1999,MartiEtAl2020}. In the baseline
analysis we focus on a full-sample MST as a summary of long-run dependence; in principle,
the same construction could be repeated on rolling windows to study how the topology
evolves over time. For a given sample, we first estimate the $N \times N$ correlation
matrix of returns
\begin{equation}
  \hat{\rho}_{ij}
  = \frac{
      \sum_{t=1}^{T^\star} (r_{i,t} - \bar{r}_i)(r_{j,t} - \bar{r}_j)
    }{
      \sqrt{\sum_{t=1}^{T^\star} (r_{i,t} - \bar{r}_i)^2}
      \sqrt{\sum_{t=1}^{T^\star} (r_{j,t} - \bar{r}_j)^2}
    },
  \qquad i,j = 1,\dots,N,
\end{equation}
where $T^\star$ is the number of observations in the estimation window and $\bar{r}_i$ is
the sample mean of $r_{i,t}$ on that window. From $\hat{\rho}_{ij}$ we define a
correlation-based distance
\begin{equation}
  d_{ij}
  = \sqrt{2(1 - \hat{\rho}_{ij})},
  \qquad d_{ij} \in [0,2],
\end{equation}
which satisfies the properties of a metric when $\hat{\rho}$ is positive semidefinite.

The complete weighted graph with nodes $\{1,\dots,N\}$ and edge weights $\{d_{ij}\}$ is
then reduced to a tree by computing the minimum spanning tree (MST) \citep{Mantegna1999}
using NetworkX's \texttt{minimum\_spanning\_tree} function, which implements Kruskal's
algorithm. The MST $\mathcal{T}$ is the connected, acyclic
subgraph that spans all $N$ nodes and minimizes the total distance
\begin{equation}
  \mathcal{T}
  = \arg\min_{\text{trees } \mathcal{G}}
    \sum_{(i,j) \in \mathcal{G}} d_{ij}.
\end{equation}
By construction, $\mathcal{T}$ contains $N-1$ edges and captures the backbone
connectivity structure of the dependence network, retaining the strongest pairwise links, by
discarding the $\binom{N}{2}-(N-1)$ weaker or redundant connections. Importantly, MST
is a network topology tool, not a statistical filter: it does not test or correct for
the significance of individual correlations, but rather extracts the minimum-weight
spanning structure of the graph. Correlations are computed with a minimum of 30 valid
joint observations; pairs below this threshold are excluded from the graph.%
\footnote{Alternative sparse network representations such as PMFG and TMFG incorporate
planarity or topological constraints and could be used as robustness checks; we adopt
MST for its computational tractability and established use in the financial network
literature \citep{Mantegna1999,MartiEtAl2020}.}

We adopt MST because correlation-based full networks constructed without multiple testing
correction may exhibit inflated edge density. MST mitigates this concern structurally by
retaining only $N-1$ edges; VAR--DY connectedness is immune to this issue entirely, as
it does not rely on individual correlation significance tests. Full correlation networks
are therefore used for visualization and exploratory analysis only, while inferential
conclusions draw on MST and VAR--DY.

From $\mathcal{T}$ we compute standard network descriptors (degree, betweenness, clustering
around hubs) and identify economically interpretable clusters (e.g.\ groups of banks or
insurers that are tightly linked in the MST).

\subsection{Country-level clustering and community detection}\label{subsec:clustering}

To examine whether national banking systems share similar systemic behaviour, we apply
community detection to the country-level TCI correlation network. For each country $c$,
we define within-country connectedness as the average share of forecast error variance
of institutions in $c$ attributable to shocks from other institutions in the same
country, computed for each rolling window $w$. The resulting time series
$\{\mathrm{TCI}_{c,w}\}$ are then summarized into a country-level correlation matrix,
and a spring-layout network is constructed from the top pairwise correlations.

Community structure in this network is estimated using the Louvain algorithm
\citep{Blondel2008}, implemented via the \texttt{python-louvain} library with fixed seed
$= 7$ to ensure reproducibility. The Louvain method maximizes a modularity objective and
does not require the number of communities $K$ to be specified in advance. This approach
groups countries that share similar domestic connectedness dynamics over the full sample,
providing an economically interpretable partition of the European banking landscape into
geographic community blocks.

\subsection{Connectedness regime identification}\label{subsec:regimes}

While the country-level clustering examines cross-sectional community structure, a
complementary question is whether the \emph{time-series} of rolling connectedness
estimates partitions naturally into distinct systemic regimes. To answer this, we apply
$k$-means clustering to a set of five summary statistics computed from the rolling GFEVD
outputs for each window $w$:

\begin{enumerate}
  \item $\mathrm{TCI}_w$:system-wide level of connectedness.
  \item $\sigma(\mathrm{NET})_w$:cross-sectional standard deviation of net spillovers,
    capturing the dispersion of institutional roles.
  \item $\max(\mathrm{NET})_w$:maximum single-institution net spillover,
    measuring the influence of the dominant transmitter.
  \item $s_5(\mathrm{TO})_w$:share of total TO connectedness concentrated in the
    top-5 emitting institutions, a Herfindahl-type concentration index.
  \item $\gamma(\mathrm{NET})_w$:skewness of the net spillover cross-section,
    indicating directional asymmetry in the network.
\end{enumerate}

All five features are standardized to zero mean and unit variance before clustering. The
optimal number of regimes $K^*$ is selected by maximising the mean silhouette coefficient
over the range $K \in \{2, 3, 4, 5, 6\}$. Final cluster centres are estimated with
$k$-means using 100 random restarts (seed $= 42$) to avoid local minima, and regimes are
labelled in ascending order of mean TCI so that Regime 1 always denotes the lower
connectedness state.

\subsection{Frequency-domain connectedness}\label{subsec:freq_conn}

Following \citet{BarunikKrehlik2018}, we decompose the aggregate time-domain
connectedness into frequency-band-specific contributions. The starting point is the
spectral representation of the VAR moving-average expansion. Define the
\emph{spectral transfer function}
\begin{equation}\label{eq:transfer}
  \Psi(e^{-i\omega})
  = \sum_{h=0}^{\infty} \Phi_h\, e^{-i\omega h},
  \qquad \omega \in [-\pi, \pi],
\end{equation}
where $\Phi_h$ are the MA coefficient matrices from Section~\ref{subsec:var_gfevd},
so that $\mathbf{y}_t = \sum_{h=0}^\infty \Phi_h \boldsymbol{\varepsilon}_{t-h}$.
The spectral density matrix of $\mathbf{y}_t$ is
\begin{equation}\label{eq:spectral_density}
  f(\omega)
  = \frac{1}{2\pi}\,\Psi(e^{-i\omega})\,\Sigma\,\Psi^*(e^{-i\omega}),
\end{equation}
where $\Psi^*(e^{-i\omega}) = \overline{\Psi(e^{-i\omega})}^{\!\top}$ denotes the
conjugate transpose, and $\mathrm{Var}(y_{i,t}) = \int_{-\pi}^{\pi} f_{ii}(\omega)\,d\omega$
by Parseval's theorem.

The generalized variance share from institution $j$ to institution $i$
\emph{at frequency} $\omega$ is
\begin{equation}\label{eq:theta_freq}
  \theta^g_{ij}(\omega)
  = \frac{
      \sigma_{jj}^{-1}\,
      \bigl|[\Psi(e^{-i\omega})\,\Sigma]_{ij}\bigr|^{2}
    }{
      [\Psi(e^{-i\omega})\,\Sigma\,\Psi^*(e^{-i\omega})]_{ii}
    },
  \qquad i,j = 1,\dots,N,
\end{equation}
the frequency-domain analogue of $\theta^g_{ij}(H)$ in
Section~\ref{subsec:var_gfevd}: it measures the contribution of institution~$j$'s
innovations to the spectral power of institution~$i$'s forecast error at cycle
length $2\pi/\omega$, normalized by the total spectral power of $i$ at that frequency.

For a frequency band $d = [\omega_1,\omega_2]\subseteq[-\pi,\pi]$, the band-specific
generalized variance share is obtained by integrating the frequency-specific
contributions weighted by the spectral power of institution~$i$:
\begin{equation}\label{eq:theta_band}
  \theta^g_{ij}(d)
  = \frac{
      \displaystyle\int_{d}
        \sigma_{jj}^{-1}\,
        \bigl|[\Psi(e^{-i\omega})\,\Sigma]_{ij}\bigr|^{2}
      \,d\omega
    }{
      \displaystyle\int_{-\pi}^{\pi}
        [\Psi(e^{-i\omega})\,\Sigma\,\Psi^*(e^{-i\omega})]_{ii}
      \,d\omega
    }.
\end{equation}
Equation~\eqref{eq:theta_band} follows from substituting~\eqref{eq:theta_freq} into
$\int_d \theta^g_{ij}(\omega)\,f_{ii}(\omega)\,d\omega / \mathrm{Var}(y_{i,t})$;
the spectral density factor $[\Psi\Sigma\Psi^*]_{ii}$ cancels in the numerator.
As in Section~\ref{subsec:var_gfevd}, shares are row-normalized:
\begin{equation}\label{eq:theta_band_norm}
  \tilde\theta^g_{ij}(d)
  = \frac{\theta^g_{ij}(d)}{\displaystyle\sum_{j=1}^{N}\theta^g_{ij}([-\pi,\pi])},
  \qquad
  \sum_{j=1}^{N}\tilde\theta^g_{ij}([-\pi,\pi]) = 1,
\end{equation}
and the band-specific Total Connectedness Index is
\begin{equation}\label{eq:tci_band}
  \mathrm{TCI}(d)
  = \frac{1}{N}
    \sum_{i=1}^{N}\sum_{\substack{j=1\\j\neq i}}^{N}
    \tilde\theta^g_{ij}(d).
\end{equation}
By construction, $\mathrm{TCI}(d_S) + \mathrm{TCI}(d_M) + \mathrm{TCI}(d_L) =
\mathrm{TCI}(H)$ in the large-$H$ limit, so the decomposition is exhaustive and
without residual.

We consider three economically motivated bands for weekly data:
\begin{itemize}
  \item \textbf{Short-run} ($d_S = [\pi/2,\,\pi]$, cycles of 2--4 weeks): captures
    fast-moving transmission typical of market stress and contagion episodes.
  \item \textbf{Medium-run} ($d_M = [\pi/4,\,\pi/2]$, cycles of 4--8 weeks): captures
    adjustment dynamics following initial shocks.
  \item \textbf{Long-run} ($d_L = [0,\pi/4]$, cycles\,$>\,$8 weeks): captures structural
    co-movement driven by common business-cycle and macroprudential exposures.
\end{itemize}
For each rolling window~$w$,
equations~\eqref{eq:transfer}--\eqref{eq:tci_band} are evaluated with the
window-specific estimates $(\Phi_{h,w}, \Sigma_w)$, yielding time-varying band-specific
TCI series $\{\mathrm{TCI}_w(d)\}_{d\in\{d_S,\,d_M,\,d_L\}}$. This decomposition
allows us to distinguish transient contagion (short-run spikes) from durable structural
integration (long-run shifts), directly addressing the question of whether COVID-19
constituted a temporary panic episode or a permanent restructuring of systemic linkages.
Frequency connectedness is computed using the \texttt{FrequencyConnectedness} function
in the \texttt{ConnectednessApproach} \textsf{R} package, with spectral resolution
\texttt{nfore}~$= 200$.

\subsection{Structural break analysis}\label{subsec:breaks}

To formally test whether the TCI series contains statistically significant breaks, rather than treating COVID-19 merely as an ex post event label, we apply the
Bai--Perron (2003) multiple-break estimation procedure and a Chow test targeted at the
pandemic onset. The Bai--Perron procedure is implemented using the \texttt{breakpoints}
function from the \texttt{strucchange} package in \textsf{R}, with minimum segment
length $h = 0.10$ (i.e.\ each regime must cover at least 10\% of the sample) and a
maximum of $K = 5$ break points. The Chow test uses 2020-03-02 as the target date,
corresponding to the week in which the COVID-19 shock first caused a marked and
sustained disruption to European financial markets.

\subsection{Network representation, portfolio design and implementation
details}\label{subsec:implementation}

For each window $w$, the normalized GFEVD matrix $\tilde{\Theta}^g_w(H)$ is interpreted
as the weighted adjacency matrix of a directed volatility-spillover network. Node $i$
corresponds to institution $i$, and the weight on the edge $i \rightarrow j$ is defined
as
\[
  w_{ij,w} = \tilde{\theta}_{ji,w}^g(H),
\]
i.e.\ the share of institution $j$'s forecast error variance attributable to shocks in
institution $i$. From this network we compute standard centrality measures (out-strength,
in-strength, eigenvector centrality, PageRank), which serve as institution-level
indicators of systemic importance and are analyzed both in levels and across clustered
regimes.

In parallel, the full-sample MST $\mathcal{T}$ derived from the correlation matrix
$\hat{\rho}_{ij}$ provides a sparse, undirected dependence network. Cluster structure in
$\mathcal{T}$ (e.g.\ communities obtained by cutting long edges or applying hierarchical
clustering on $d_{ij}$) is used to group institutions into diversified blocks. In
principle, simple portfolio rules can be constructed that (i)~overweight
low-connectedness or net-receiver institutions (based on $C_{i,w}^{\text{net}}$), and
(ii)~allocate capital according to correlation/MST-based clustering, in the spirit of
hierarchy- and network-based allocation schemes surveyed by \citet{MartiEtAl2020}. We
outline this link between network structure and portfolio design as a natural application
of the framework; implementing and backtesting such strategies in detail is left for
future work.

In practice, VAR estimation and GFEVD computation are implemented in \textsf{R} using
standard connectedness routines, while clustering, MST computation, and portfolio
experiments are implemented in \textsf{Python}. Hyperparameters such as VAR order $p$,
forecast horizon $H$, window length $W$, number of clusters $K$, and the choice of
rolling window for correlation estimation are treated as tuning parameters: they are
chosen based on information criteria, sensitivity analysis, and economic
interpretability, and their final values are reported in the empirical section together
with robustness checks.

All empirical steps, from data preprocessing and weekly aggregation to VAR--DY
estimation, DCC--GARCH modelling, and network construction, are implemented in a
reproducible \textsf{R}/\textsf{Python} pipeline.
